\section*{Broader Impact Statement}

The dominant paradigm for deploying machine learning systems might be characterized as ``train, deploy, and pray''. Collect data, train a model, tune hyperparameters until empirical performance seems acceptable, deploy, and hope that test-time behavior remains within tolerable bounds. When failures occur, the response is typically reactive, patching the system after harm has been observed. This paradigm persists not because practitioners prefer it, but because principled alternatives have been lacking.

This work contributes to a shift from safety-through-patching toward safety-by-design. Rather than discovering failure modes empirically and addressing them post-hoc, conformal policy control allows practitioners to specify acceptable risk levels before deployment and obtain guarantees that these levels will be respected. The hyperparameter tuning phase becomes about improving performance within safety constraints, rather than searching for a configuration that might satisfy unspecified feasibility requirements.

If such methods become widespread, the implications for where and how machine learning is deployed could be substantial. High-stakes domains that have resisted ML adoption due to liability concerns or regulatory barriers, such as clinical decision support, autonomous systems, and financial services, may become more accessible when developers can provide formal guarantees rather than empirical assurances. Such a change in framing would align ML deployment with the safety certification practices already standard in aviation, pharmaceuticals, and other regulated industries.

